% ═══════════════════════════════════════════════════════════════
% Section 3: From the Hourglass to an Addressing Model
% (label kept as sec:hourglass so earlier cross-references resolve)
% ═══════════════════════════════════════════════════════════════

\section{From the Hourglass to an Addressing Model}
\label{sec:hourglass}

Versions up to v1.5 of this paper presented an 11-band \emph{hourglass}: seven neural
bands (N7 neocortex down to N1 spinal cord), one interface band (I0, the electrode-tissue
boundary) and three synthetic bands (S1--S3), read as a stack. The model borrowed from the
OSI reference model~\cite{zimmermann1980osi} and from the Internet's narrow
waist~\cite{deering1998hourglass}, and made three claims: that I0 is a natural security
chokepoint through which all signals pass; that every component of a BCI maps to exactly
one band, so that no attack surface falls outside the model; and that behaviour becomes
less deterministic ascending the neural bands. This revision withdraws all three and
replaces the stack with addresses. This section records what changed and why.

\subsection{Why the Stack Was Withdrawn}

\paragraph{The numbering carries no flow.}
If the bands were a stack, neural pathways should mostly join numerically adjacent bands.
Recomputing each pathway's origin and target bands from the structure table gives 65
links between distinct N codes, of which 18 (27.7\%) join adjacent codes. Over all 5,040
possible numberings of the seven codes the mean is 28.6\%, and 60\% of numberings do at
least as well as the published one. The order is no better than chance.

\paragraph{The thalamus is not a gate.}
The v1.5 text used thalamocortical relay~\cite{sherman2006thalamus} to support
signal-plausibility checks, and earlier project documents described the thalamus as the
brain's firewall. The thalamus is the principal route by which sensory input, cerebellar
output and basal ganglia output reach neocortex, and transmission through it is
state-dependent~\cite{sherman2016thalamuscentral}. It is not a security control:
brainstem neuromodulatory systems also reach cortex by extra-thalamic
routes~\cite{edlow2012ascendingarousal}, motor thalamus integrates cortical, cerebellar and
basal ganglia input~\cite{boschbouju2013motorthalamus}, and olfaction is not thalamus-free,
since piriform cortex projects to the mediodorsal thalamus~\cite{courtiol2015olfactorythalamus}.

\paragraph{I0 is a boundary, not a chokepoint.}
The electrode-tissue boundary is where a recording or stimulating device transduces between
ionic and electronic current, and it remains an address. It is not a point through which
all attacks pass: of the \LiveTechniques{} techniques in the live catalogue, 110 carry no
I0 tag and 53 are tagged only at device codes, among them consumer-sensor side channels
that never touch an electrode.

\paragraph{The model was not complete.}
Nerves, roots, ganglia and sense organs had no address, so the completeness claim did not
hold. A periphery code (PER) is proposed for them and is not yet integrated.

\paragraph{Determinacy was untested.}
The per-band determinacy classes and ``QI'' ranges had no defined estimator. They are now
recorded as untested hypotheses (\texttt{datalake/qif-hypotheses.json}), proposed for the
earlier band memberships.

\subsection{Addressing Version 2}

The seven N codes are retained as identifiers. Each structure has one address, set by the
adult territory that holds its principal neurons; circuits such as the basal ganglia and
the limbic system are functional tags a structure can carry, not addresses.
Table~\ref{tab:bands} gives the codes in use. \textbf{The number carries no severity,
processing or traversal meaning.}

\begin{table}[H]
\centering
\caption{Address codes in use (addressing version 2). Listed by code, not by rank.}
\label{tab:bands}
\small
\begin{tabularx}{\textwidth}{@{}l l X@{}}
\toprule
\textbf{Code} & \textbf{Name} & \textbf{Membership rule} \\
\midrule
N7 & Neocortex       & Isocortex, including agranular motor areas \\
N6 & Allocortex      & Allocortex, mesocortex and pallial nuclei, including basolateral amygdala (v1.5: ``Limbic System'') \\
N5 & Cerebral nuclei & Subpallial nuclei: striatum, pallidum, basal forebrain, central and medial amygdala (v1.5: ``Basal Ganglia'') \\
N4 & Diencephalon    & Thalamus, subthalamus, hypothalamus, epithalamus \\
N3 & Cerebellum      & Cerebellar cortex and nuclei \\
N2 & Brainstem       & Midbrain, pons and medulla, including substantia nigra and ventral tegmental area \\
N1 & Spinal cord     & Spinal cord parenchyma; roots and ganglia are outside it \\
\midrule
I0 & Neural interface & Electrode-tissue boundary where biological signals become electronic \\
\midrule
S1 & Analog front-end   & Amplification, filtering, ADC/DAC conversion \\
S2 & Digital processing & Decoding algorithms, classification, digital signal processing \\
S3 & Application        & Clinical software, user interface, data storage, network \\
\bottomrule
\end{tabularx}
\end{table}

Three structures changed address. The substantia nigra and ventral tegmental area moved
from N5 to N2, because they lie in the midbrain, and the subthalamic nucleus moved from N5
to N4. All three keep their basal ganglia or reward tags. Subthalamic deep brain
stimulation, which v1.5 placed at N5, is therefore an N4 example. Because codes keep their
identifiers while membership changes, the atlas carries \texttt{addressing\_version: 2},
which should be cited with any band reference. Pathway bands are now computed from the
structure table by a script with a check mode, which the build does not yet run.

Fourteen statements in project documents drew security conclusions from the gate analogy,
chiefly that a deeper implant bypasses fewer protective layers. They are marked under
review, neither deleted nor asserted. The NISS formula has no term for implant depth: all
\LiveNissScored{} stored scores reproduce exactly from their six metric values, although
that does not show whether depth reasoning influenced the choice of any metric level.

\subsection{Routes (Proposed)}

Signal flow, which the stack could not carry, is moved to a proposed route view, staged and
not integrated. A route is an ordered list of synaptic hops, each with its sign,
transmitter and own evidence. A source check of all 90 hops in the 17 staged routes found
47 supported by a source a checker read, 42 partly supported and 1 contradicted for
primates; 3 hops are recorded as demonstrated in humans and 85 as inferred by homology from
other species. The routes are suitable for tagging what a device touches and for teaching.
They are not a basis for claims about what an attack can reach in a person.

\paragraph{Effect on this paper.}
The band annotations in the snapshot analysis (Sections~\ref{sec:tara}--\ref{sec:nic}) use
the version 1 memberships. Those analyses are not recomputed here; a reader using a band
count from them should treat it as a version 1 figure.
